\documentclass[11pt]{article}
\usepackage{fontspec,amsmath,amssymb,amsthm,hyperref}
\usepackage[margin=1in]{geometry}
\setlength{\emergencystretch}{3em}
\title{Exact scaling of coordinate smoothness constants}
\author{LangevinLean mathematical exposition}
\date{}
\begin{document}\maketitle
\noindent Source module: \texttt{AverageLMC/SmoothnessScaling.lean}. Compile this proof document with XeLaTeX or LuaLaTeX. The accompanying HTML pairs the mathematical argument with the exact Lean statements and records the scope of the formalization.


This module proves that the literal supremum constants commute with nonnegative rescaling of the potential. In particular the barred average-coordinate constant of the SGLD average potential f/n is exactly the full-potential constant divided by n. The proof does not replace either quantity by an arbitrary upper bound.



For a differentiable potential with Lipschitz gradient, nonnegative multiplication commutes with the literal coordinate Taylor suprema:

\[c\ge0\quad\Longrightarrow\quad M_j(cf)=cM_j(f),\qquad M_{\rm av}(cf)=cM_{\rm av}(f).\]

In particular the SGLD barred constant is exactly the constant of the average potential:

\[n\ge1\quad\Longrightarrow\quad \bar M_{\rm av}=M_{\rm av}(f/n)=\frac{M_{\rm av}(f)}{n}.\]
\section{Rescale the actual gradient and coordinate Taylor quotients}


The Frechet derivative of c times f is c times the derivative of f. The Riesz representation defining the gradient commutes with this real scalar, so the true gradient is c times the original gradient. Substitution in the exact coordinate Taylor quotient factors out c. Consequently the entire quotient set of the rescaled potential is the image of the original set under multiplication by c.



\noindent Lean declarations: \texttt{gradient\_const\_mul\_potential}, \texttt{coordinateTaylorQuotient\_const\_mul}, \texttt{coordinateTaylorQuotients\_const\_mul}.

\section{Transport the least upper bound through nonnegative multiplication}


The Lipschitz-gradient assumption supplies an upper bound on the original Taylor-quotient set, and that set is nonempty. Its supremum therefore is a genuine least upper bound. Nonnegative multiplication transports that least upper bound to the image set, including scalar zero. This proves exact homogeneity of every coordinate supremum. A finite sum and division by the dimension prove the same identity for M\_av.



\noindent Lean declarations: \texttt{coordinateSmoothness\_const\_mul}, \texttt{averageSmoothness\_const\_mul}.

\section{Identify the paper's barred average-potential constant}


Choose the scalar to be the reciprocal of the finite component count. It is nonnegative, and real division is multiplication by this reciprocal. The proved homogeneity yields M\_av(x ↦ f(x)/n)=M\_av(f)/n. Thus the source-facing SGLD theorem can state its discretization term with the literal constant of the average potential.



\noindent Lean declarations: \texttt{averageSmoothness\_div\_nat}.

\end{document}
